\documentclass{amsart}
% For dealing with references we use the comment environment
\usepackage{verbatim}
\newenvironment{reference}{\comment}{\endcomment}
%\newenvironment{reference}{}{}
\newenvironment{slogan}{\comment}{\endcomment}
\newenvironment{history}{\comment}{\endcomment}

% For commutative diagrams we use Xy-pic
\usepackage[all]{xy}

% We use 2cell for 2-commutative diagrams.
\xyoption{2cell}
\UseAllTwocells

% We use multicol for the list of chapters between chapters
\usepackage{multicol}

% This is generally recommended for better output
\usepackage{lmodern}
\usepackage[T1]{fontenc}

% For cross-file-references

% Package for hypertext links:
\usepackage{hyperref}

% For any local file, say "hello.tex" you want to link to please

% Theorem environments.
%
\theoremstyle{plain}
\newtheorem{theorem}[subsection]{Theorem}
\newtheorem{proposition}[subsection]{Proposition}
\newtheorem{lemma}[subsection]{Lemma}

\theoremstyle{definition}
\newtheorem{definition}[subsection]{Definition}
\newtheorem{example}[subsection]{Example}
\newtheorem{exercise}[subsection]{Exercise}
\newtheorem{situation}[subsection]{Situation}

\theoremstyle{remark}
\newtheorem{remark}[subsection]{Remark}
\newtheorem{remarks}[subsection]{Remarks}

\numberwithin{equation}{subsection}

% Macros
%
\def\lim{\mathop{\mathrm{lim}}\nolimits}
\def\colim{\mathop{\mathrm{colim}}\nolimits}
\def\Spec{\mathop{\mathrm{Spec}}}
\def\Hom{\mathop{\mathrm{Hom}}\nolimits}
\def\Ext{\mathop{\mathrm{Ext}}\nolimits}
\def\SheafHom{\mathop{\mathcal{H}\!\mathit{om}}\nolimits}
\def\SheafExt{\mathop{\mathcal{E}\!\mathit{xt}}\nolimits}
\def\Sch{\mathit{Sch}}
\def\Mor{\mathop{\mathrm{Mor}}\nolimits}
\def\Ob{\mathop{\mathrm{Ob}}\nolimits}
\def\Sh{\mathop{\mathit{Sh}}\nolimits}
\def\NL{\mathop{N\!L}\nolimits}
\def\CH{\mathop{\mathrm{CH}}\nolimits}
\def\proetale{{pro\text{-}\acute{e}tale}}
\def\etale{{\acute{e}tale}}
\def\QCoh{\mathit{QCoh}}
\def\Ker{\mathop{\mathrm{Ker}}}
\def\Im{\mathop{\mathrm{Im}}}
\def\Coker{\mathop{\mathrm{Coker}}}
\def\Coim{\mathop{\mathrm{Coim}}}

% Boxtimes
%
\DeclareMathSymbol{\boxtimes}{\mathbin}{AMSa}{"02}

%
% Macros for moduli stacks/spaces
%
\def\QCohstack{\mathcal{QC}\!\mathit{oh}}
\def\Cohstack{\mathcal{C}\!\mathit{oh}}
\def\Spacesstack{\mathcal{S}\!\mathit{paces}}
\def\Quotfunctor{\mathrm{Quot}}
\def\Hilbfunctor{\mathrm{Hilb}}
\def\Curvesstack{\mathcal{C}\!\mathit{urves}}
\def\Polarizedstack{\mathcal{P}\!\mathit{olarized}}
\def\Complexesstack{\mathcal{C}\!\mathit{omplexes}}
% \Pic is the operator that assigns to X its picard group, usage \Pic(X)
% \Picardstack_{X/B} denotes the Picard stack of X over B
% \Picardfunctor_{X/B} denotes the Picard functor of X over B
\def\Pic{\mathop{\mathrm{Pic}}\nolimits}
\def\Picardstack{\mathcal{P}\!\mathit{ic}}
\def\Picardfunctor{\mathrm{Pic}}
\def\Deformationcategory{\mathcal{D}\!\mathit{ef}}

\setlength{\emergencystretch}{3em}
\begin{document}
\title[Stacks Project, Tag 0FEW]{456 --- Rational coherent K theory agrees with Chow homology for schemes embedded in regular quasi-compact schemes}
\maketitle
\noindent Copyright (C) 2005--2025 Johan de Jong. Stacks Project authors. Source: \href{https://stacks.math.columbia.edu/tag/0FEW}{Tag 0FEW}.

\noindent Editorial difficulty note (not author text). Difficulty H2: The localized Chern character must be proved triangular for the support-dimension filtration, not simply assumed to be an isomorphism. The proof tracks supports of arbitrary coherent classes and determines the leading coefficient of every integral support by reduction to a generic regular-sequence embedding, yielding inverse maps on all graded pieces..

\noindent Editorial provenance note (not author text). History: excerpt from revision \texttt{a0444\allowbreak 6e57e\allowbreak c1fbc\allowbreak 252a8\allowbreak 71afc\allowbreak ec775\allowbreak 2fb28\allowbreak 07b14}, chow.tex, lines 13190--13281. Reference presentation was adapted; original-author context and core support blocks were retained or added. The mathematical wording is unchanged. Licensed under GNU FDL 1.2 or later; no invariant sections or cover texts. See ../licenses/COPYING. Context blocks reproduce author definitions and supporting results; they are not additional counted problems.

\section{Chow groups and K-groups revisited}
\label{section-chow-and-K-II}

\noindent
This section is the continuation of Section \href{https://stacks.math.columbia.edu/tag/0FDQ}{section chow and K (Tag 0FDQ)}.
Let $(S, \delta)$ be as in Situation \ref{situation-setup}.
Let $X$ be locally of finite type over $S$. The K-group
$K'_0(X) = K_0(\textit{Coh}(X))$ of coherent sheaves on $X$
has a canonical increasing filtration
$$
F_kK'_0(X) =
\Im\Big(K_0(\textit{Coh}_{\leq k}(X)) \to K_0(\textit{Coh}(X)\Big)
$$
This is called the filtration by dimension of supports. Observe that
$$
\text{gr}_k K'_0(X) \subset K'_0(X)/F_{k - 1}K'_0(X) =
K_0(\textit{Coh}(X)/\textit{Coh}_{\leq k - 1}(X))
$$
where the equality holds
by Homology, Lemma \href{https://stacks.math.columbia.edu/tag/02MX}{lemma serre subcategory K groups (Tag 02MX)}.
The discussion in Remark \href{https://stacks.math.columbia.edu/tag/02SD}{remark good cases K A (Tag 02SD)} shows
that there are canonical maps
$$
\CH_k(X) \longrightarrow \text{gr}_k K'_0(X)
$$
defined by sending the class of an integral closed subscheme
$Z \subset X$ of $\delta$-dimension $k$ to the class of
$[\mathcal{O}_Z]$ on the right hand side.



\begin{definition}
\label{definition-bivariant-class}
\begin{reference}
Similar to \href{https://stacks.math.columbia.edu/bibliography}{Bibliography: F (Definition 17.1)}
\end{reference}
Let $(S, \delta)$ be as in Situation \ref{situation-setup}.
Let $f : X \to Y$ be a morphism of schemes locally of finite type over $S$.
Let $p \in \mathbf{Z}$.
A {\it bivariant class $c$ of degree $p$ for $f$} is given by a rule
which assigns to every locally of finite type morphism $Y' \to Y$
and every $k$ a map
$$
c \cap - : \CH_k(Y') \longrightarrow \CH_{k - p}(X')
$$
where $X' = Y' \times_Y X$, satisfying the following conditions
\begin{enumerate}
\item if $Y'' \to Y'$ is a proper, then
$c \cap (Y'' \to Y')_*\alpha'' = (X'' \to X')_*(c \cap \alpha'')$
for all $\alpha''$ on $Y''$ where $X'' = Y'' \times_Y X$,
\item if $Y'' \to Y'$ is flat locally of finite type of
fixed relative dimension, then
$c \cap (Y'' \to Y')^*\alpha' = (X'' \to X')^*(c \cap \alpha')$
for all $\alpha'$ on $Y'$, and
\item if $(\mathcal{L}', s', i' : D' \to Y')$ is as in
Definition \ref{definition-gysin-homomorphism}
with pullback $(\mathcal{N}', t', j' : E' \to X')$ to $X'$,
then we have $c \cap (i')^*\alpha' = (j')^*(c \cap \alpha')$
for all $\alpha'$ on $Y'$.
\end{enumerate}
The collection of all bivariant classes of degree $p$ for $f$ is
denoted $A^p(X \to Y)$.
\end{definition}

\begin{situation}
\label{situation-setup}
Here $S$ is a locally Noetherian, and universally catenary scheme.
Moreover, we assume $S$ is endowed with a dimension function
$\delta : S \longrightarrow \mathbf{Z}$.
\end{situation}

\begin{definition}
\label{definition-gysin-homomorphism}
Let $(S, \delta)$ be as in Situation \ref{situation-setup}.
Let $X$ be locally of finite type over $S$.
Let $(\mathcal{L}, s)$ be a pair consisting of an invertible
sheaf and a global section $s \in \Gamma(X, \mathcal{L})$.
Let $D = Z(s)$ be the zero scheme of $s$, and
denote $i : D \to X$ the closed immersion.
We define, for every integer $k$, a {\it Gysin homomorphism}
$$
i^* : Z_{k + 1}(X) \to \CH_k(D).
$$
by the following rules:
\begin{enumerate}
\item Given an integral closed subscheme $W \subset X$ with
$\dim_\delta(W) = k + 1$ we define
\begin{enumerate}
\item if $W \not \subset D$, then $i^*[W] = [D \cap W]_k$ as a
$k$-cycle on $D$, and
\item if $W \subset D$, then
$i^*[W] = i'_*(c_1(\mathcal{L}|_W) \cap [W])$,
where $i' : W \to D$ is the induced closed immersion.
\end{enumerate}
\item For a general $(k + 1)$-cycle $\alpha = \sum n_j[W_j]$
we set
$$
i^*\alpha = \sum n_j i^*[W_j]
$$
\item If $D$ is an effective Cartier divisor, then we denote
$D \cdot \alpha = i_*i^*\alpha$ the pushforward of the class $i^*\alpha$
to a class on $X$.
\end{enumerate}
\end{definition}

\begin{proposition}
\label{proposition-K-tensor-Q}
Let $(S, \delta)$ be as in Situation \ref{situation-setup}. Assume given a
closed immersion $X \to Y$ of schemes locally of finite type over $S$
with $Y$ regular and quasi-compact. Then the composition
$$
K'_0(X) \to
K_0(D_{X, perf}(\mathcal{O}_Y)) \to
A^*(X \to Y) \otimes \mathbf{Q} \to
\CH_*(X) \otimes \mathbf{Q}
$$
of the map $\mathcal{F} \mapsto \mathcal{F}[0]$ from
Remark \href{https://stacks.math.columbia.edu/tag/0FEQ}{remark perf Z regular (Tag 0FEQ)}, the map $ch(X \to Y, -)$ from
Remark \href{https://stacks.math.columbia.edu/tag/0FET}{remark localized chern character K (Tag 0FET)}, and
the map $c \mapsto c \cap [Y]$ induces an isomorphism
$$
K'_0(X) \otimes \mathbf{Q}
\longrightarrow
\CH_*(X) \otimes \mathbf{Q}
$$
which depends on the choice of $Y$. Moreover, the canonical map
$$
\CH_k(X) \otimes \mathbf{Q}
\longrightarrow
\text{gr}_k K'_0(X) \otimes \mathbf{Q}
$$
(see above) is an isomorphism of $\mathbf{Q}$-vector spaces for all
$k \in \mathbf{Z}$.
\end{proposition}

\begin{proof}
Since $Y$ is regular, the construction in
Remark \href{https://stacks.math.columbia.edu/tag/0FEQ}{remark perf Z regular (Tag 0FEQ)} applies.
Since $Y$ is quasi-compact, the construction in
Remark \href{https://stacks.math.columbia.edu/tag/0FET}{remark localized chern character K (Tag 0FET)} applies.
We have that $Y$ is locally equidimensional
(Lemma \href{https://stacks.math.columbia.edu/tag/0FE3}{lemma locally equidimensional (Tag 0FE3)}) and
thus the ``fundamental cycle'' $[Y]$ is defined
as an element of $\CH_*(Y)$, see Remark \href{https://stacks.math.columbia.edu/tag/0FE4}{remark fundamental class (Tag 0FE4)}.
Combining this with the map $\CH_k(X) \to \text{gr}_kK'_0(X)$
constructed above we see that it suffices to prove
\begin{enumerate}
\item If $\mathcal{F}$ is a coherent $\mathcal{O}_X$-module
whose support has $\delta$-dimension $\leq k$, then
the composition above sends $[\mathcal{F}]$ into
$\bigoplus_{k' \leq k} \CH_{k'}(X) \otimes \mathbf{Q}$.
\item If $Z \subset X$ is an integral closed subscheme
of $\delta$-dimension $k$, then the composition above
sends $[\mathcal{O}_Z]$ to an element whose degree $k$
part is the class of $[Z]$ in $\CH_k(X) \otimes \mathbf{Q}$.
\end{enumerate}
Namely, if this holds, then our maps induce maps
$\text{gr}_kK'_0(X) \otimes \mathbf{Q} \to CH_k(X) \otimes \mathbf{Q}$
which are inverse to the canonical maps 
$\CH_k(X) \otimes \mathbf{Q} \to \text{gr}_k K'_0(X) \otimes \mathbf{Q}$
given above the proposition.

\medskip\noindent
Given a coherent $\mathcal{O}_X$-module $\mathcal{F}$
the composition above sends $[\mathcal{F}]$ to
$$
ch(X \to Y, \mathcal{F}[0]) \cap [Y] \in \CH_*(X) \otimes \mathbf{Q}
$$
If $\mathcal{F}$ is (set theoretically) supported on a closed subscheme
$Z \subset X$, then we have
$$
ch(X \to Y, \mathcal{F}[0]) = (Z \to X)_* \circ ch(Z \to Y, \mathcal{F}[0])
$$
by Lemma \href{https://stacks.math.columbia.edu/tag/0FB7}{lemma loc chern shrink Z (Tag 0FB7)}. We conclude that in this
case we end up in the image of $\CH_*(Z) \to \CH_*(X)$. Hence
we get condition (1).

\medskip\noindent
Let $Z \subset X$ be an integral closed subscheme of $\delta$-dimension $k$.
The composition above sends $[\mathcal{O}_Z]$ to the element
$$
ch(X \to Y, \mathcal{O}_Z[0]) \cap [Y] =
(Z \to X)_* ch(Z \to Y, \mathcal{O}_Z[0]) \cap [Y]
$$
by the same argument as above.
Thus it suffices to prove that the degree $k$ part of
$ch(Z \to Y, \mathcal{O}_Z[0]) \cap [Y] \in
\CH_*(Z) \otimes \mathbf{Q}$ is $[Z]$.
Since $\CH_k(Z) = \mathbf{Z}$, in order
to prove this we may replace $Y$ by an open neighbourhood of the
generic point $\xi$ of $Z$. Since the maximal ideal of the regular
local ring $\mathcal{O}_{X, \xi}$ is generated by a
regular sequence (Algebra, Lemma \href{https://stacks.math.columbia.edu/tag/00NQ}{lemma regular ring CM (Tag 00NQ)})
we may assume the ideal of $Z$ is generated by a regular sequence, see
Divisors, Lemma \href{https://stacks.math.columbia.edu/tag/063I}{lemma Noetherian scheme regular ideal (Tag 063I)}.
Thus we deduce the result from Lemma \href{https://stacks.math.columbia.edu/tag/0FEH}{lemma actual computation (Tag 0FEH)}.
\end{proof}
\end{document}
